\def\BibFileName{template.bib}                  % name of BibLaTeX file with all the 
\def\FigurePath{figures}                        % subdirectory for the figure files 

%%%%%%%%%%%%%%%%%%%%%%%% LaTeX CLASS AND THEME %%%%%%%%%%%%%%%%%%%%%%%%
\documentclass[xcolor=dvipsnames,aspectratio=169]{beamer}
\useoutertheme[left]{sidebar}
\useinnertheme{rounded}
\usefonttheme[onlymath]{serif}                  % to use regular math font

%%%%%%%%%%%%%%%%%%%%%%%%%%%%% JH colors %%%%%%%%%%%%%%%%%%%%%%%%%%%%%%%
%%%% following colors are defined based on: https://brand.jhu.edu/color
\definecolor{HeritageBlue}{RGB}{0,45,114}       % heritage blue
\definecolor{SpiritBlue}{RGB}{114,172,229}      % spirit blue

\definecolor{JHUPastle}{RGB}{0,155,119}         % pastle
\definecolor{JHUDeepBlue}{RGB}{0,114,206}       % deep blue

\definecolor{JHURed}{RGB}{166,25,46}            % red
\definecolor{JHUGreen}{RGB}{40,97,64}           % green

\definecolor{JHUViolet}{RGB}{81,40,79}          % violet
\definecolor{JHUPurple}{RGB}{161,90,149}        % purple
%%%%%%%%%%%%%%%%%%%%%%%%%%%%% JH COLORS %%%%%%%%%%%%%%%%%%%%%%%%%%%%%%%

%%%%%%%%%%%%%%%%%%%%% COLOR PALETTE SETTINGS %%%%%%%%%%%%%%%%%%%%%%%%%%
%%%% JHU color is used to set the palette for the beamer template

\setbeamercolor{palette primary}{bg=HeritageBlue,fg=white}
\setbeamercolor{palette secondary}{bg=HeritageBlue,fg=white}
\setbeamercolor{palette tertiary}{bg=HeritageBlue,fg=white}
\setbeamercolor{palette quaternary}{bg=HeritageBlue,fg=white}
\setbeamercolor{sidebar}{bg=White}
\setbeamercolor{palette sidebar primary}{bg=white,fg=HeritageBlue}
\setbeamercolor{palette sidebar secondary}{bg=white,fg=HeritageBlue}
\setbeamercolor{palette sidebar tertiary}{bg=white,fg=HeritageBlue}
\setbeamercolor{palette sidebar quaternary}{bg=white,fg=HeritageBlue}

\setbeamercolor{titlelike}{bg=HeritageBlue,fg=White}
\setbeamercolor{author}{fg=HeritageBlue}
\setbeamercolor{institute}{fg=HeritageBlue}
\setbeamercolor{date}{fg=HeritageBlue}

\setbeamercolor{frametitle}{bg=HeritageBlue,fg=White}
\setbeamercolor{frametitle continuation}{bg=HeritageBlue,fg=White}

\setbeamercolor{section in toc}{fg=HeritageBlue}        % TOC sections
\setbeamercolor{structure}{fg=HeritageBlue}             % itemize, enumerate, etc

\setbeamercolor{title in sidebar}{fg=HeritageBlue}
\setbeamercolor{author in sidebar}{fg=SpiritBlue}

\setbeamercolor{block title}{bg=HeritageBlue,fg=White}
\setbeamercolor{block body}{bg=HeritageBlue!10,fg=Black} 
\setbeamercolor{block title alerted}{bg=JHURed,fg=White}
\setbeamercolor{block body alerted}{bg=JHURed!10,fg=Black} 
\setbeamercolor{block title example}{bg=JHUGreen,fg=White}
\setbeamercolor{block body example}{bg=JHUGreen!10,fg=Black} 

\setbeamercolor{footnote}{bg=White,fg=HeritageBlue}
\setbeamercolor{footnote mark}{bg=White,fg=HeritageBlue}

\setbeamercolor{alerted text}{fg=JHURed}

%%%%%%%%%%%%%%%%%%% END COLOR PALETTE SETTINGS %%%%%%%%%%%%%%%%%%%%%%%%

%%%%%%%%%%%%%%%%%%%%%%%%%%%% FONT SETTINGS %%%%%%%%%%%%%%%%%%%%%%%%%%%%
%%% size of the fonts in different environments of the beamer template
\setbeamerfont{title}{size=\Large}
\setbeamerfont{subtitle}{size=\large}
\setbeamerfont{author}{size=\normalsize}
\setbeamerfont{institute}{size=\scriptsize}
\setbeamerfont{frametitle}{size=\Large}

%%%%%%%%%%%%%%%%%%%%%%%%%% END FONT SETTINGS %%%%%%%%%%%%%%%%%%%%%%%%%%

%%%%%%%%%%%%%%% FRAME TITLE, SIDEBAR, FOOTER SETTINGS %%%%%%%%%%%%%%%%%
\makeatletter
    \logo{\includegraphics[height=1cm]{JHU_shield.pdf}}
    \setlength{\beamer@headheight}{1cm}
\makeatother

\makeatletter
\setlength\beamer@sidebarwidth{0.75in}
  \setbeamertemplate{sidebar \beamer@sidebarside}%{sidebar theme}
  {
    \beamer@tempdim=\beamer@sidebarwidth%
    \advance\beamer@tempdim by -6pt%
    \insertverticalnavigation{\beamer@sidebarwidth}%
    \vfill
    \ifx\beamer@sidebarside\beamer@lefttext%
    \else%
      \usebeamercolor{normal text}%
      \llap{\usebeamertemplate***{navigation symbols}\hskip0.1cm}%
      \vskip2pt%
    \fi%
  }%
\makeatother

\setbeamertemplate{navigation symbols}
{\tiny\insertframenumber/\inserttotalframenumber}

%%%%%%%%%%%%% END FRAME TITLE, SIDEBAR, FOOTER SETTINGS %%%%%%%%%%%%%%%

%%%%%%%%%%%%%%%%%%%%%%%%%%% TOC SETTINGS %%%%%%%%%%%%%%%%%%%%%%%%%%%%%%
%%%% the next block of commands puts the table of contents at the 
%%%% beginning of each section and highlights the current section
\AtBeginSection[]
{
  \begin{frame}
    \frametitle{Table of Contents}
    \tableofcontents[currentsection]
  \end{frame}
}
%%%%%%%%%%%%%%%%%%%%%%%%% END TOC SETTINGS %%%%%%%%%%%%%%%%%%%%%%%%%%%%
%%%%%%%%%%%%%%%%%%%%%% END LaTeX CLASS AND THEME %%%%%%%%%%%%%%%%%%%%%%

%%%%%%%%%%%%%%%%%%%%% LaTeX PACKAGE AND SETTINGS %%%%%%%%%%%%%%%%%%%%%%
\usepackage{amsfonts,amssymb,amsmath,amsthm,autobreak, cancel,dsfont,mathtools,mathbbol,mathrsfs,siunitx,upgreek}


\usepackage[ruled]{algorithm2e}             % to manage algorithm environment
\usepackage[american]{babel}                % for different language typography

\usepackage[style=authoryear,autocite=footnote,backend=biber]{biblatex}

\usepackage{blindtext}

\usepackage{caption,subcaption}             % to manage captions
\usepackage{csquotes,epigraph,varwidth}     % for managing quotes
\usepackage{float}                          % to manage floating environment
\usepackage[T1]{fontenc}                    % for font encoding
\usepackage[bottom]{footmisc}               % footnote environment management
\usepackage{graphicx,wrapfig}               % to manage images
\usepackage{geometry}                       % to manage margins and others
\usepackage{hyperref}                       % for hyperlinks
\usepackage{listings}                       % to include codes

%%% table related packages
\usepackage{booktabs,longtable,makecell,multicol,multirow,tabularx,xltabular}

\usepackage[rightcaption]{sidecap}          % for sideway captions
\usepackage{tikz}                           % package for drawing

\hypersetup{colorlinks,linkcolor=,urlcolor=JHUPurple,citecolor=JHUPurple}
\addbibresource{\BibFileName}
\graphicspath{{\FigurePath/}}

%%%%%%%%%%%%%%%%%%% END LaTeX PACKAGE AND SETTINGS %%%%%%%%%%%%%%%%%%%%

%%%%%%%%%%%%%%%%%%%%%%%%% LaTeX MACROS %%%%%%%%%%%%%%%%%%%%%%%%%%%%%%%%

\newcommand{\BLUE}{\textcolor{JHUDeepBlue}}
\newcommand{\PURPLE}{\textcolor{JHUPurple}}
\newcommand{\RED}{\textcolor{Red}}
\newcommand{\GREEN}{\textcolor{Green}}

%%%%%%%%%%%%%%%%%%%%%%% END LaTeX MACROS %%%%%%%%%%%%%%%%%%%%%%%%%%%%%%
%%%%%%%%%%%%%%%%%%%%%%%%%%%%%%%%%%%%%%%%%%%%%%%%%%%%%%%%%%%%%%%%%%%%%%%
%%%%%%%%%%%%%%%%%%%%%%%%%%%%%%%%%%%%%%%%%%%%%%%%%%%%%%%%%%%%%%%%%%%%%%%


%%%%%%%%%%%%%%%%%%%%%%%%%%%%% DOCUMENT INFORMATION %%%%%%%%%%%%%%%%%%%%%%%%%%%%%
\title[About Beamer]
{Beamer template for Johns Hopkins University}
%\subtitle{Template made on customized theme}
\author[Doe]{John Doe}

\institute[JHU] {
    Affiliated position \\
    Name of department or institute \\
    \includegraphics[width=0.35\textwidth, trim={0.6in 0.35in 0.6in 0.3in},clip]{JHU_logo.pdf}
}

\date[VLC 2021]{Very Large Conference, April 2021}
%%%%%%%%%%%%%%%%%%%%%%%%%%% END DOCUMENT INFORMATION %%%%%%%%%%%%%%%%%%%%%%%%%%%

\begin{document}

\begingroup
\makeatletter
\setbeamertemplate{navigation symbols}{}
\setlength{\hoffset}{-.4\beamer@sidebarwidth}
\makeatother
\begin{frame}[plain]
    \titlepage
\end{frame}
\endgroup

%%%%%%%%%%%%%%%%%%%%%%%%%%%%%%%%%%%%%%%%%%%%%%%%%%%%%%%%%%%%%%%%%%%%%%%%%%%%%%%%
\section{First project}
\subsection{Introduction}
\begin{frame}{Highlighting text}
    
    In this slide, some important text will be \alert{highlighted} with JHU color because it's important. Please, don't \alert{abuse} it.
    
    \begin{block}{Remark}
        Sample text
    \end{block}
    
    \begin{alertblock}{Important theorem}
        Sample text in the red box
    \end{alertblock}
    
    \begin{examples}
        Sample text in the green box. The title of the block is ``Examples''.
    \end{examples}
\end{frame}

%%%%%%%%%%%%%%%%%%%%%%%%%%%%%%%%%%%%%%%%%%%%%%%%%%%%%%%%%%%%%%%%%%%%%%%%%%%%%%%%
\subsection{Environment tests}
\begin{frame}{Enumerate and itemize test}
    \textbf{Enumerate environment with classic colors}
    \begin{enumerate}
        \item Hello, this is the \RED{first} point
        %
        \item This is my \GREEN{second} point
    \end{enumerate}
    
    \vspace{0.5in}
    
    \textbf{Itemize environment with JHU colors}
    \begin{itemize}
        \item This is a separate \BLUE{itemize} test 
        %
        \item Second point for the \PURPLE{separate} test
    \end{itemize}
\end{frame}
%%%%%%%%%%%%%%%%%%%%%%%%%%%%%%%%%%%%%%%%%%%%%%%%%%%%%%%%%%%%%%%%%%%%%%%%%%%%%%%%
%%%%%%%%%%%%%%%%%%%%%%%%%%%%%%%%%%%%%%%%%%%%%%%%%%%%%%%%%%%%%%%%%%%%%%%%%%%%%%%%
\section{Second project}
\subsection{Testing math environment}
\begin{frame}{Integration}
    Integration is also known as anti-derivative which represents the area under the curve. For indefinite integral, we should add a constant $C$ to the anti-derivative expressions.
    \begin{equation*}
        \int \exp(x) \ dx = \exp(x) + C
    \end{equation*}
    
    You can \textbf{\BLUE{boldface colors}} to emphasize even more\footnote{Hyperlinks will appear as \url{https://overleaf.com}}.
   
\end{frame}


%%%%%%%%%%%%%%%%%%%%%%%%%%%%%%%%%%%%%%%%%%%%%%%%%%%%%%%%%%%%%%%%%%%%%%%%%%%%%%%%
\subsection{Testing overlays}
\begin{frame}{More Lists}
    \begin{enumerate}[I]
    \item<1-> Point A\footnotemark
    \item<2-> Point B
    \begin{itemize}
    \item<3-> part 1
    \item<4-> part 2
    \end{itemize}
    \item<5-> Point C
    \item<6-> Point D
    \end{enumerate}
    
    \footnotetext{Second footnote testing here}
\end{frame}
%%%%%%%%%%%%%%%%%%%%%%%%%%%%%%%%%%%%%%%%%%%%%%%%%%%%%%%%%%%%%%%%%%%%%%%%%%%%%%%%
\subsection{Overlaying math environment}
\begin{frame}{Maths Blocks}
    \begin{theorem}<1->[Pythagoras] 
        $ a^2 + b^2 = c^2$
    \end{theorem}
    \begin{corollary}<2->
        $ x + y = y + x  $
    \end{corollary}
    \begin{proof}<3->
        $\omega +\phi = \epsilon $
    \end{proof}
\end{frame}
%%%%%%%%%%%%%%%%%%%%%%%%%%%%%%%%%%%%%%%%%%%%%%%%%%%%%%%%%%%%%%%%%%%%%%%%%%%%%%%%
%%%%%%%%%%%%%%%%%%%%%%%%%%%%%%%%%%%%%%%%%%%%%%%%%%%%%%%%%%%%%%%%%%%%%%%%%%%%%%%%
\section{Third project}
\subsection{References}
\begin{frame}{References}
    \blindtext \footcite{einstein}
\end{frame}
%%%%%%%%%%%%%%%%%%%%%%%%%%%%%%%%%%%%%%%%%%%%%%%%%%%%%%%%%%%%%%%%%%%%%%%%%%%%%%%%

\end{document}